\section{Conclusion}
In this work, we demonstrate how to learn scalable, convolutional cells from data that transfer to multiple image classification tasks. The learned architecture is quite flexible as it may be scaled in terms of computational cost and parameters to easily address a variety of problems. In all cases, the accuracy of the resulting model exceeds all human-designed models -- ranging from models designed for mobile applications to computationally-heavy models designed to achieve the most accurate results.

The key insight in our approach is to design a search space that decouples the complexity of an architecture from the depth of a network. This resulting search space permits identifying good architectures on a small dataset (i.e., CIFAR-10) and transferring the learned architecture to image classifications across a range of data and computational scales.

The resulting architectures approach or exceed state-of-the-art performance in both CIFAR-10 and ImageNet datasets with less computational demand than human-designed architectures \cite{szegedy2016rethinking,BatchNorm,zhang2016polynet}. The ImageNet results are particularly important because many state-of-the-art computer vision problems (e.g.,  object detection \cite{huang2016speed}, face detection \cite{schroff2015facenet}, image localization \cite{weyand2016planet}) derive image features or architectures from ImageNet classification models. For instance, we find that image features obtained from ImageNet used in combination with the Faster-RCNN framework achieves state-of-the-art object detection results. Finally, we demonstrate that we can use the resulting learned architecture to perform ImageNet classification with reduced computational budgets that outperform streamlined architectures targeted to mobile and embedded platforms \cite{howard2017mobilenets, shufflenet}.

%In this work, we introduce a new search space for Neural Architecture Search that decouples the complexity of the architecture from the depth of the convolutional neural network. This search space permites identifying good architectures on CIFAR-10 and transferring these architectures to ImageNet. The resulting architectures exhibit state-of-the-art predictive performance but contain fewer parameters and require fewer computational operations.
%With our new search space, Neural Architecture Search tends to find good architectures that have fewer trainable parameters and are relatively shallow compared to state-of-the-art models.
%On ImageNet, our best architectures achieve a great tradeoff between predictive accuracy versus the number of trainable parameters and number of floating point operations. 

%Our results have strong implications for transfer learning and meta-learning as this is the first work to demonstrate state-of-the-art results using meta-learning on a large scale problem.
%This work also highlights that learned elements of network architectures, beyond model weights, can transfer across datasets.
%The degree to which meta-learning approaches may learn architectures that exceed human-invented architectures in domains outside of image classification remains an open question for exploration.

%An interesting future direction is to explore how automated learning may exceed human-invented architectures in domains outside of image classification.
%These results suggest that automated learning may learn architectures that surpass

%in domains outside of image classification may exceed human-invented architectures \cite{object-detection, embedding, unsupervised-learning}.

%The fact that the architectures are transferable between datasets such as CIFAR-10 and ImageNet also indicates that there are other elements of architectures beyond weight-sharing that can be transferred across datasets. 

%\todo{Emphasize comment: 5: Maybe stress the fact that the learned architecture can be scaled up and down in terms of computational cost and parameters to easily provide a variety of tradeoffs, and that in all cases, the accuracy for these tradeoffs exceeds all known human-designed models at the equivalent levels of computational cost, ranging from very lightweight models designed for mobile applications all the way up to computational heavy models designed to get the most accurate results.}
